\section{Results}
\label{sec:results}

\subsection{Controlled H200 comparison}

% Generated by scripts/generate_controlled_tables.py; do not edit.
\begin{table}[t]
\caption{Matched held-out loss by seed. Lower is better. Each value is the final step-3,000 evaluation over 1,048,576 held-out tokens.}
\label{tab:h200-controlled}
\centering
\small
\begin{tabular}{lrrr}
\toprule
Seed & GQA & W/R/V & W/R/V $-$ GQA \\
\midrule
42 & 4.683425 & \textbf{4.662531} & $-0.020894$ \\
43 & 4.717673 & \textbf{4.686053} & $-0.031620$ \\
44 & 4.708008 & \textbf{4.704711} & $-0.003297$ \\
\midrule
Mean & 4.703035 & \textbf{4.684432} & $-0.018604$ \\
Sample SD & 0.017657 & 0.021137 & 0.014300 \\
\bottomrule
\end{tabular}
\end{table}

\begin{table}[t]
\caption{Training throughput on one H200. Each seed contributes its post-warmup median after evaluation/checkpoint stalls are excluded; the table reports the mean and sample SD across seeds.}
\label{tab:h200-throughput}
\centering
\small
\begin{tabular}{lrr}
\toprule
Mechanism & Mean train tok/s $\uparrow$ & Sample SD \\
\midrule
GQA & \textbf{125,021} & 100 \\
Contextual W/R/V & 121,448 & 254 \\
\bottomrule
\end{tabular}
\end{table}


W/R/V is lower in loss on all three seeds. The mean paired difference is $-0.018604$, approximately 0.40\% of the GQA mean loss. The two-sided 95\% Student-$t$ interval is $[-0.0541,0.0169]$. Because it includes zero and contains only three pairs, we describe the ordering as consistent in this experiment rather than statistically established.

W/R/V is 3,573 tokens/s slower on average, a 2.86\% throughput penalty. Quality and throughput therefore move in opposite directions: the controlled result is a modest quality improvement at modest training-speed cost, not a Pareto improvement.

\subsection{Historical architecture path}

Historical single-seed H100 runs motivated the controlled comparison but are not matched to the H200 protocol. Their key completed losses are 4.895427 for collision-normalized fixed-state WVR, 4.892026 after adding one W/R/V layer, 4.973798 for a strongly lexical eight-layer W/R/V condition, 4.852761 for an eight-layer contextual hybrid, and 4.706118 for the ten-layer contextual condition with RMS normalization. These values document the development path only. They are not combined with Table~\ref{tab:h200-controlled}, and no causal credit is assigned from their joint changes.

Short proxies were useful failure screens but unreliable rankers: several favorable lexical-forge and multi-timescale proxy results did not improve at full budget. The final controlled model therefore removes the lexical write residual from the evaluated W/R/V attention path. Tied lexical objects and deterministic micro-experts remain in the shared feed-forward architecture of both GQA and W/R/V.

Two H200 seed-42 ablations are reported descriptively and are not pooled with the three-seed comparison. Enabling the lexical write residual gives 4.660939 loss versus 4.662531 with the residual disabled, a difference of $-0.001592$, while reducing median throughput from 121,479 to 119,121 tokens/s ($-1.94\%$). Disabling per-head R/W RMS normalization gives 4.795811 loss, $+0.133280$ relative to the lexical-off normalized condition, at 123,257 tokens/s. The lexical-residual run uses an earlier source commit, and both comparisons have one seed; these observations motivate matched ablations but do not isolate causal effects.

\subsection{Fixed-state negative controls}

The historical fixed-state experiments remain informative negative controls. A separate 1B-token causal-convolution run cannot be ranked against the 100M-token conditions because its budget differs by a factor of ten. Archived associative-recall and induction probes report zero accuracy for evaluated attention-free checkpoints, but the available GQA probe does not document an identical protocol; no quantitative positive-control comparison is claimed.

FP32 unit tests reproduce full-sequence W/R/V outputs under token-by-token incremental execution at tolerance $10^{-5}$. The retained W/V tensors grow with context, so these tests support causal implementation equivalence rather than a fixed-state serving claim.